\section{Introduction}\label{sec:intro}

Clearing the denominators in a rational recurrence produces polynomial
relations that remain meaningful at zero coordinates. Their solutions may
carry nilpotent structure that is invisible in the recurrence's original
domain. Even when the possible zero patterns are understood, determining
the lengths of the corresponding local algebras is a separate intersection
problem. This distinction is particularly concrete for the exchange
recurrence
\begin{equation}\label{eq:recurrence}
 x_{i-1}x_{i+1}=1+x_i^k,\qquad k\geq3\ \text{odd}.
\end{equation}
Imposing a cyclic clock gives a finite algebra. At an alternating-zero
point its local embedding dimension grows with the clock, so a description
of the support on a two-dimensional cluster variety cannot by itself
determine the retained infinitesimal structure.

The support in question is classical in this setting. Beyer and Muller
classify the deep points of coefficient-free rank-two cluster algebras;
their alternating pattern is $(0,\alpha,0,\alpha^{-1})$, with the relevant
root condition, and the shifted phase
\cite[Section~4, especially Corollary~4.6 in arXiv v1]{BM}.
We use this attribution while deriving the elementary finite-clock
restriction directly. Benito, Faber, Mourtada, and Schober classify the
ambient rank-two singularities; their Jacobian description implies
smoothness for the positive rank-two exponents over $\C$
\cite[Section~5, Theorem~5.2.3 in arXiv v2]{BFMS}.
The smooth ambient surface and the unsaturated cyclic relation algebra
studied here are different schemes. A further related, but distinct,
direction is the birational dynamics of cluster mutation maps studied by
Grigorev, Kalidindi, Quintero Santander, and Roeder \cite{GKQR}.
For equal exponents their mutation composition corresponds to the square
of $F_k(x,y)=(y,(1+y^k)/x)$. We do not use a Picard-action or entropy
calculation to infer a boundary length.

Our result determines the local thickness of the entire alternating-zero
stratum, simultaneously for every odd $k\geq3$ and every clock at which it
occurs. The local equations reduce to products of coordinates and gradient
components of one formal potential. Expanding their intersection
multiplicity over all factor choices leaves path and cycle gradient
colengths. Nonsingular Hessians account for the simplest terms, but they
do not settle the problem: paths of odd length have a one-dimensional
kernel, and cycles of length divisible by four have a two-dimensional
kernel. The cancelled path case has leading effective degree $2k$ rather
than $k$. The cycle case requires both a mixed quartic coefficient and a
pure-axis coefficient. Both calculations have additional terms at $k=3$.

The resulting weights depend only on path length modulo four. This makes
the local lengths amenable to a labeled cyclic-word enumeration and yields
a rational multiplicity-generating function. The rational series is a
consequence of the same local intersection calculation, not a new
dynamical realization. The original rational recurrence excludes these
zero-coordinate points from its admissible torus domain, and other
boundary supports remain outside the theorem.

\Cref{sec:main} states the formula and fixes the small-clock conventions.
\Cref{sec:formal} proves the formal reduction and the required nonreduced
intersection additivity. The path and cycle calculations occupy
\cref{sec:paths,sec:cycles}. We derive the rational series in
\cref{sec:series}, and distinguish the finite computational controls from
the proof and the unresolved boundary problem in \cref{sec:limits}.
